\section{Introduction}

Entity matching (EM) is a core step in data integration projects: given two collections of records, decide which pairs refer to the same real-world entity \cite{christophides2020overview,christophides2015webdata}. Modern EM systems increasingly rely on supervised models, including transformer-based matchers \cite{mudgal2018deep,brunner2020transformer,li2020ditto,barlaug2021survey}, but these models need labeled pairs before they can be deployed in a new domain \cite{konda2016magellan,peeters2024wdcproducts}. Creating such training data is expensive because useful labels must include not only obvious matches and non-matches, but also boundary cases: same-family products, compatible components, near-duplicate publication titles, and other pairs that are easy to confuse \cite{christen2012data,peeters2024wdcproducts}.

Large language models (LLMs) make it possible to revisit this bottleneck. Prior work has shown that LLMs can directly classify entity pairs in zero-shot settings and can be fine-tuned for EM \cite{peeters2025llmem,steiner2025finetuning}. In this paper, we study a different, more operational question. If a strong LLM can label pairs, is it still necessary to manually label training data for every new EM task, or can the training-data construction step be automated?

We answer this question by studying whether an LLM-based system can \emph{build an entire training dataset by selecting pairs from the two source tables and labeling those pairs}. The goal is not to replace the downstream matcher with an LLM at inference time. Instead, the LLM is used as a data construction tool: it labels selected examples once, after which the resulting training set can be used to train a conventional downstream matcher such as Ditto or a smaller local model.

This distinction matters because a useful training-set construction system has to solve two coupled problems. First, it must select which candidate pairs should be included in the final training set. A random or nearest-neighbor-only strategy may spend most of the annotation budget on easy negatives, while missing positives or hard non-matches. Second, it must label the selected pairs. The LLM component is responsible for assigning match/non-match labels, but the downstream value of the selected-and-labeled training set depends just as much on the selection component as on the labeling component.

The main question of the paper is therefore practical:

%\begin{quote}
%Can an LLM-based training-set construction system replace manual EM training-data labeling by selecting and labeling training sets that match or beat the official training sets of standard benchmarks?
%\end{quote}

\begin{quote}
Can a LLM combined with an active learning-based pair selection strategy label training sets that have the same quality as the training sets that are part of widely-used entity matching benchmarks?
\end{quote}

Training set quality = Performance of downstream system~\cite{FelixDataQualityandML2025}.

Our conclusion: 
1. Current LLMs combined with an active learning-based pair selection strategy can label training sets having same or a better quality as the training sets that are part of widely-used entity matching benchmarks.
2. This means that a large part of the training data labeling process can be automated, resulting in substantially reduced costs. 
3. If in addition to the LLM also human effort is invested in labeling training sets, this effort can be directed towards deciding tricky corner-cases, for instance cases where different LLMs disagree or the LLM outputs a low confidence in its decision.


We answer this question empirically by constructing selected-and-labeled training sets for multiple EM benchmarks, training downstream matchers on these selected-and-labeled sets, and comparing their performance to matchers trained on the official benchmark training splits. Preliminary results show that selected-and-labeled training sets can perform equally or better than official training sets on several benchmarks. The second part of the paper then asks why this happens. We compare selected-and-labeled and official training sets by pair overlap, entity coverage, class balance, similarity and difficulty distributions, and qualitative categories of selected examples.

The central finding is that the result is not only a labeling result. Selection matters. A simple similarity-search strategy can spend much of its budget on repetitive or easy negatives. An active-learning strategy, especially one that uses the downstream matcher in the loop, selects more positives and more decision-boundary examples. This changes both downstream F1 and labeling cost. We therefore treat the workflow as a combined \emph{selection-and-labeling} problem rather than as a pure LLM classification problem.

The paper is organized around five research questions.

\begin{description}
  \item[RQ1] Can selected-and-labeled training sets match or outperform official benchmark training sets for downstream EM?
  \item[RQ2] Which example-selection strategies and training-set sizes are most effective?
  \item[RQ3] How expensive is LLM-based training-set construction, and how does cost differ across selection strategies?
  \item[RQ4] Do the findings depend on one proprietary LLM labeler, or do they transfer to other strong LLMs, including open-weight models?
  \item[RQ5] How do selected-and-labeled training sets differ from official training sets at the pair, entity, label, and difficulty levels, and which differences explain downstream performance gains or losses?
\end{description}

The intended contributions are:

\begin{itemize}
  \item an LLM-based system for selecting and labeling EM training data that combines candidate selection with LLM labeling;
  \item an empirical comparison of selected-and-labeled versus official training sets across standard EM benchmarks;
  \item evidence that active-learning-based example selection produces more useful training data than simpler similarity-based selection;
  \item a cost analysis showing that targeted selection can reduce the amount of LLM labeling needed;
  \item a robustness analysis plan across different LLM labelers, including an open-weight labeling model for settings where data cannot be sent to a hosted API;
  \item a detailed analysis methodology for explaining what selected-and-labeled training sets contain and how they differ from official splits.
\end{itemize}
